\documentclass[11pt,letterpaper]{article}
\usepackage[T1]{fontenc}
\usepackage[scaled]{helvet}
\renewcommand{\familydefault}{\sfdefault}
\usepackage{amsmath,amssymb}
\usepackage[margin=0.8in,bottom=0.9in]{geometry}
\usepackage[most]{tcolorbox}
\usepackage{tikz}
\usetikzlibrary{calc}
\usepackage{enumitem}
\usepackage{array}
\usepackage{needspace}
\usepackage{fancyhdr}
\usepackage{pifont}

% ---- Palette "F. Bright & Playful" ----
\definecolor{fmtitle}{HTML}{3D2C8D}
\definecolor{fmkey}{HTML}{0096C7}
\definecolor{fmwatch}{HTML}{2BB673}
\definecolor{fmturn}{HTML}{F4A100}
\definecolor{fmmistake}{HTML}{EF476F}
\definecolor{fmnotice}{HTML}{9B5DE5}
\definecolor{fmgray}{HTML}{6B6B6B}

\pagestyle{fancy}
\fancyhf{}
\renewcommand{\headrulewidth}{0pt}
\fancyfoot[C]{\footnotesize\color{fmgray}\textcopyright\ 2026 FangMath. All rights reserved.}
\fancyfoot[R]{\footnotesize\color{fmgray}\thepage}
\setlength{\parindent}{0pt}
\setlength{\parskip}{3pt}

\newcommand{\badge}[1]{\colorbox{fmtitle}{\color{white}\bfseries\small\ #1\ }}
\newcommand{\leveltag}[2]{\colorbox{#1}{\color{white}\bfseries\scriptsize\ #2\ }}
\newcommand{\LE}{\leveltag{fmwatch}{EASY}}
\newcommand{\LM}{\leveltag{fmturn}{MEDIUM}}
\newcommand{\LH}{\leveltag{fmmistake}{HARD}}

\newcommand{\sectionbar}[2]{\needspace{6cm}\par\vspace{4mm}
  \begin{tcolorbox}[enhanced,colback=fmtitle,colframe=fmtitle,arc=2mm,
    left=4mm,top=1.5mm,bottom=1.5mm]
  \color{white}\bfseries\large #1\hfill{\normalsize\mdseries #2}
  \end{tcolorbox}}

\newtcolorbox{keyidea}{enhanced,breakable,colback=white,colframe=fmkey,boxrule=1pt,arc=2mm,
  title={Key Idea},coltitle=fmkey,fonttitle=\bfseries,
  attach boxed title to top left={xshift=4mm,yshift=-2mm},
  boxed title style={colback=white,colframe=fmkey,boxrule=0.9pt}}
\newtcolorbox{watchme}{enhanced,colback=white,colframe=fmwatch,boxrule=1pt,arc=2mm,
  title={Watch Me --- Worked Example},coltitle=fmwatch,fonttitle=\bfseries,
  attach boxed title to top left={xshift=4mm,yshift=-2mm},
  boxed title style={colback=white,colframe=fmwatch,boxrule=0.9pt}}
\newtcolorbox{mistakebox}{enhanced,breakable,colback=white,colframe=fmmistake,boxrule=1pt,arc=2mm,
  title={Common Mistake},coltitle=fmmistake,fonttitle=\bfseries,
  attach boxed title to top left={xshift=4mm,yshift=-2mm},
  boxed title style={colback=white,colframe=fmmistake,boxrule=0.9pt}}
\newtcolorbox{yourturn}{enhanced,breakable,colback=white,colframe=fmturn,boxrule=1pt,arc=2mm,
  title={Your Turn},coltitle=fmturn,fonttitle=\bfseries,
  attach boxed title to top left={xshift=4mm,yshift=-2mm},
  boxed title style={colback=white,colframe=fmturn,boxrule=0.9pt}}
\newtcolorbox{answerbox}[1]{enhanced,breakable,colback=white,colframe=fmwatch,boxrule=1pt,arc=2mm,
  title={#1},coltitle=fmwatch,fonttitle=\bfseries,
  attach boxed title to top left={xshift=4mm,yshift=-2mm},
  boxed title style={colback=white,colframe=fmwatch,boxrule=0.9pt}}

\newcommand{\work}[1]{\par\vspace{#1}}
\setlist[enumerate]{itemsep=2pt,topsep=3pt}

% ---------------- extra boxes ----------------
\newtcolorbox{notebox}{enhanced,breakable,colback=white,colframe=fmnotice,boxrule=1pt,arc=2mm,
  title={Notice},coltitle=fmnotice,fonttitle=\bfseries,
  attach boxed title to top left={xshift=4mm,yshift=-2mm},
  boxed title style={colback=white,colframe=fmnotice,boxrule=0.9pt}}
\newtcolorbox{exbox}{enhanced,colback=white,colframe=fmwatch,boxrule=0.8pt,arc=2mm,
  title={Example},coltitle=fmwatch,fonttitle=\bfseries\small,
  attach boxed title to top left={xshift=4mm,yshift=-2mm},
  boxed title style={colback=white,colframe=fmwatch,boxrule=0.8pt},fontupper=\small}
\newcommand{\topicbar}[1]{\needspace{4.5cm}\par\vspace{5mm}
  \begin{tcolorbox}[enhanced,colback=fmtitle,colframe=fmtitle,arc=2mm,left=4mm,top=1.5mm,bottom=1.5mm]
  \color{white}\bfseries\large #1
  \end{tcolorbox}}
\usepackage{multicol}

% ---------------- practice layout ----------------
\newcommand{\workline}{\par\vspace{7mm}{\color{gray!55}\rule{\linewidth}{0.4pt}}}
\newcommand{\worklines}[1]{\foreach \i in {1,...,#1}{\workline}}
\newcommand{\prob}[3]{\needspace{#3}\par\vspace{3mm}{\color{fmturn}\bfseries #1.}\ #2\par}
\newcommand{\ansline}{\par\vspace{4mm}\hfill\textbf{Answer:}\ \rule{5cm}{0.4pt}\par}
\newcommand{\keyitem}[3]{\needspace{3.2cm}\par\vspace{2mm}
  \begin{tcolorbox}[enhanced,colback=white,colframe=fmwatch!60,boxrule=0.6pt,arc=1.5mm,left=3mm,top=1.5mm,bottom=1.5mm]
  {\color{fmwatch}\bfseries #1.}\quad\small #2\par\vspace{1mm}
  \normalsize\textbf{Answer: #3}
  \end{tcolorbox}}


\usetikzlibrary{arrows.meta}
\providecommand{\iskey}{0}
% fill-in blank: line for students, pink answer in the key
\newcommand{\fb}[2][2.6cm]{\ifnum\iskey=1 \underline{\makebox[#1][c]{\color{fmmistake}\bfseries #2}}\else\underline{\makebox[#1]{}}\fi}
% answer line for practice problems
\newcommand{\pans}[2][3.2cm]{\ifnum\iskey=1 \underline{\makebox[#1][c]{\color{fmmistake}\bfseries #2}}\else\rule{#1}{0.4pt}\fi}
\newcommand{\vid}[1]{{\color{fmgray}\small\textit{\ding{43} Video: #1}}}
\newcommand{\pt}[2]{\fill (#1) circle (2.2pt) node[above=1pt,font=\small]{$#2$};}

\usetikzlibrary{decorations.pathreplacing}
% three points on a segment with lengths written above the pieces: #1,#2,#3 letters; #4,#5 piece labels; #6 whole label (below)
\newcommand{\threept}[6]{\begin{tikzpicture}[x=1cm]\draw[thick] (0,0)--(3.6,0);
\fill (0,0) circle(2.2pt) node[below=2pt,font=\small]{$#1$} (1.4,0) circle(2.2pt) node[below=2pt,font=\small]{$#2$} (3.6,0) circle(2.2pt) node[below=2pt,font=\small]{$#3$};
\node[above,font=\small] at (0.7,0.02){#4};\node[above,font=\small] at (2.5,0.02){#5};
\if\relax\detokenize{#6}\relax\else\draw[gray,decorate,decoration={brace,amplitude=4pt,mirror}] (0,-0.45)--(3.6,-0.45) node[midway,below=4pt,font=\small]{#6};\fi
\end{tikzpicture}}
\newcommand{\sdp}[3]{\begin{minipage}[t][2.4cm]{0.48\linewidth}{\color{fmturn}\bfseries #1.}\ #2\par\vfill\hfill \pans[3.6cm]{#3}\end{minipage}}
\newcommand{\figp}[4]{\begin{minipage}[t]{0.48\linewidth}{\color{fmturn}\bfseries #1.}\ #2\par\vspace{1mm}\centering #3\par\vspace{2mm}\raggedleft \pans[2.6cm]{#4}\end{minipage}}
\newcommand{\algp}[3]{\begin{minipage}[t][4.6cm]{0.48\linewidth}{\color{fmturn}\bfseries #1.}\ #2\par\vfill\hfill \pans[4.2cm]{#3}\end{minipage}}


% ch2: long key answers wrap inside their column instead of running into the next one
\newlength{\pw}
\renewcommand{\pans}[2][3.2cm]{\ifnum\iskey=1 \settowidth{\pw}{\bfseries #2}\ifdim\pw>#1 \parbox[b]{#1}{\raggedleft\color{fmmistake}\bfseries\small #2}\else\underline{\makebox[#1][c]{\color{fmmistake}\bfseries #2}}\fi\else\rule{#1}{0.4pt}\fi}
\begin{document}
\ifnum\iskey=1 \colorbox{fmwatch}{\color{white}\bfseries\small\ ANSWER KEY\ }\else\badge{PRACTICE}\fi\hfill{\small\color{fmgray}Geometry \textperiodcentered\ Lesson 2.1}\par\vspace{4mm}
{\color{fmtitle}\LARGE\bfseries Patterns \& Inductive Reasoning --- Practice\ifnum\iskey=1 \ --- Answer Key\fi}\par\vspace{1mm}
{\color{fmgray}A worked example comes before each type. Show your thinking.}\par\vspace{3mm}
Name: \rule{6cm}{0.4pt}\hspace{1cm}Date: \rule{3.5cm}{0.4pt}\par\vspace{2mm}

\needspace{6cm}\topicbar{Type 1 \quad Find the Next Terms}
\begin{watchme}$5,\ 11,\ 17,\ 23,\ \ldots$ \quad Check every pair: $11 - 5 = 6$, $17 - 11 = 6$, $23 - 17 = 6$. The rule is \textbf{add 6}.\par
Next two terms: $23 + 6 = 29$, \ $29 + 6 = 35$. \quad \textbf{Answer: $\boldsymbol{29,\ 35}$}\end{watchme}
\needspace{3.4cm}\noindent\sdp{1}{$4,\ 9,\ 14,\ 19,\ \ldots$ \ Next two terms?}{$24,\ 29$}\hfill\sdp{2}{$50,\ 43,\ 36,\ 29,\ \ldots$ \ Next two terms?}{$22,\ 15$}\par\vspace{2mm}
\needspace{3.4cm}\noindent\sdp{3}{$3,\ 6,\ 12,\ 24,\ \ldots$ \ Next two terms?}{$48,\ 96$}\hfill\sdp{4}{$1,\ 4,\ 9,\ 16,\ 25,\ \ldots$ \ Next two terms?}{$36,\ 49$}\par\vspace{2mm}
\needspace{3.4cm}\noindent\sdp{5}{$2,\ 3,\ 5,\ 8,\ 12,\ \ldots$ \ Next two terms?}{$17,\ 23$}\hfill\sdp{6}{$1,\ -2,\ 4,\ -8,\ \ldots$ \ Next two terms?}{$16,\ -32$}\par

\needspace{6cm}\topicbar{Type 2 \quad Find a Far-Away Term}
\begin{watchme}$7,\ 10,\ 13,\ 16,\ \ldots$ \ Find the 20th term.\par
Start at $7$ and add $3$ \textbf{nineteen} times (one fewer than the term number): $7 + 19 \cdot 3 = 7 + 57$. \quad \textbf{Answer: $\boldsymbol{64}$}\end{watchme}
\needspace{3.4cm}\noindent\sdp{7}{$2,\ 7,\ 12,\ 17,\ \ldots$ \ Find the 10th term.}{$47$}\hfill\sdp{8}{$100,\ 96,\ 92,\ 88,\ \ldots$ \ Find the 15th term.}{$44$}\par\vspace{2mm}
\needspace{3.4cm}\noindent\sdp{9}{Figure 1 has $4$ toothpicks, Figure 2 has $7$, Figure 3 has $10$. How many in Figure 25?}{$76$}\hfill\sdp{10}{\LH\ $1,\ 3,\ 6,\ 10,\ 15,\ \ldots$ \ Find the 10th term.}{$55$}\par

\needspace{6cm}\topicbar{Type 3 \quad Make a Conjecture}
\begin{watchme}Look at sums of two odd numbers: $1 + 3 = 4$, \ $5 + 7 = 12$, \ $9 + 11 = 20$.\par
Every sum is even. \quad \textbf{Conjecture: the sum of two odd numbers is even.}\end{watchme}
\needspace{3.4cm}\noindent\sdp{11}{$2 \cdot 4 = 8$, $\ 6 \cdot 2 = 12$, $\ 10 \cdot 8 = 80$. Make a conjecture about the product of two even numbers.}{the product is even}\hfill\sdp{12}{$1 + 3 = 4$, $\ 1 + 3 + 5 = 9$, $\ 1 + 3 + 5 + 7 = 16$. Conjecture about the sum of the first $n$ odd numbers.}{it equals $n^2$}\par

\needspace{6cm}\topicbar{Type 4 \quad Find a Counterexample}
\begin{watchme}Conjecture: ``If a number is divisible by $2$, it is divisible by $4$.''\par
Try $6$: it is divisible by $2$, but $6 \div 4 = 1.5$. One counterexample shows the conjecture is false. \quad \textbf{Answer: $\boldsymbol{6}$}\end{watchme}
\needspace{3.4cm}\noindent\sdp{13}{``The square of a number is always greater than the number.'' Give a counterexample.}{$\tfrac12$ (or $1$, $0$)}\hfill\sdp{14}{``Every prime number is odd.'' Give a counterexample.}{$2$}\par\vspace{2mm}
\needspace{3.4cm}\noindent\sdp{15}{``Two angles that add to $90^\circ$ are both $45^\circ$.'' Give a counterexample.}{$30^\circ$ and $60^\circ$}\hfill\sdp{16}{\LH\ ``$n^2 + n + 1$ is prime for every whole number $n$.'' Find a counterexample.}{$n = 4$: $21 = 3 \cdot 7$}\par
\end{document}
